% GENERATED FILE. Edit canonical lessons, then run npm run generate:books.
% canonical-source-hash: fnv1a64:1e6cdba041d9251a
% canonical-lessons: BN-C339-odhikar, BN-C339-dayitvo, BN-C339-svartho, BN-C339-rokkha-kora, BN-C339-bhog-kora

\chapter{Rights and Duties}
\label{ch:bn-rights-and-duties}
\hlchaptermodality{339}

\begin{chapteropening}
\textbf{By the end of this chapter:} \emph{I can say what my rights and my duties are.}

Complete the last lesson of chapter 339: I can say what my rights and my duties are.
\end{chapteropening}

\section[\texorpdfstring{\bn{অধিকার} (ôdhikār)}{অধিকার (ôdhikār)}]{\bn{অধিকার} (ôdhikār) — a right (to something)}

\label{lesson:BN-C339-odhikar}



\begin{quote}
Before the new one: say the Bengali for to check against, then the Bengali for to refute.
\end{quote}

\subsection*{You'll want to know: \bn{অধিকার}}
\textbf{\bn{অধিকার}} — \emph{ôdhikār} — \textquotedblleft{}a right (to something)\textquotedblright{}.

\begin{quote}
\bn{ভোট} \bn{দেওয়া} \bn{নাগরিকের} \bn{অধিকার।}
\end{quote}

\begin{quote}
\emph{bhoṭ deoyā nāgôriker ôdhikār.}
\end{quote}

\begin{quote}
Voting is a citizen's right.
\end{quote}

\begin{cousinweb}
\textbf{\bn{অধিকার}} (\emph{ôdhikār}) is Sanskrit \emph{adhikāra}, authority, a claim or a right (Monier-Williams). It belongs to someone: \bn{নাগরিকের} \bn{অধিকার}.
\end{cousinweb}

\begin{grammarlens}[title={What you've built}]
One more word for giving reasons and opinions.
\end{grammarlens}

\subsection*{Your turn}
\begin{itemize}
\raggedright
  \item \emph{Say it:} \emph{ôdhikār}
  \item \emph{Say it:} \emph{ôdhikār}, once more
  \item \emph{Say it:} \emph{bhoṭ deoyā nāgôriker ôdhikār.} — Voting is a citizen's right.
  \item \emph{Recall:} say the Bengali for to check against, then the Bengali for to refute, then say \emph{ôdhikār} again
\end{itemize}

\begin{tcolorbox}[breakable,colback=teal!4,colframe=teal!35!black,title={Before you move on}]
What does \bn{অধিকার} mean? (\textquotedblleft{}A right (to something)\textquotedblright{}.) Say it once more.
\end{tcolorbox}
\section[\texorpdfstring{\bn{দায়িত্ব} (dāyitvô)}{দায়িত্ব (dāyitvô)}]{\bn{দায়িত্ব} (dāyitvô) — a responsibility}

\label{lesson:BN-C339-dayitvo}



\begin{quote}
Before the new one: say the Bengali for to refute, then the Bengali for a right.
\end{quote}

\subsection*{You'll want to know: \bn{দায়িত্ব}}
\textbf{\bn{দায়িত্ব}} — \emph{dāyitvô} — \textquotedblleft{}a responsibility\textquotedblright{}.

\begin{quote}
\bn{পরিবারকে} \bn{সাহায্য} \bn{করা} \bn{আমার} \bn{দায়িত্ব।}
\end{quote}

\begin{quote}
\emph{pôribārke sāhājjo kôrā āmār dāyitvô.}
\end{quote}

\begin{quote}
It is my responsibility to help my family.
\end{quote}

\begin{cousinweb}
\textbf{\bn{দায়িত্ব}} (\emph{dāyitvô}) is a learned formation: Sanskrit \emph{dāyin}, which Monier-Williams gives as giving and also as owing or having to pay, with the abstract ending \emph{-tva}: the state of owing something to others. The word itself is not in the dictionary.
\end{cousinweb}

\begin{grammarlens}[title={What you've built}]
One more word for giving reasons and opinions.
\end{grammarlens}

\subsection*{Your turn}
\begin{itemize}
\raggedright
  \item \emph{Say it:} \emph{dāyitvô}
  \item \emph{Say it:} \emph{dāyitvô}, once more
  \item \emph{Say it:} \emph{pôribārke sāhājjo kôrā āmār dāyitvô.} — It is my responsibility to help my family.
  \item \emph{Recall:} say the Bengali for to refute, then the Bengali for a right, then say \emph{dāyitvô} again
\end{itemize}

\begin{tcolorbox}[breakable,colback=teal!4,colframe=teal!35!black,title={Before you move on}]
What does \bn{দায়িত্ব} mean? (\textquotedblleft{}A responsibility\textquotedblright{}.) Say it once more.
\end{tcolorbox}
\section[\texorpdfstring{\bn{স্বার্থ} (svārthô)}{স্বার্থ (svārthô)}]{\bn{স্বার্থ} (svārthô) — self-interest; one's own interest}

\label{lesson:BN-C339-svartho}



\begin{quote}
Before the new one: say the Bengali for a right, then the Bengali for a responsibility.
\end{quote}

\subsection*{You'll want to know: \bn{স্বার্থ}}
\textbf{\bn{স্বার্থ}} — \emph{svārthô} — \textquotedblleft{}self-interest; one's own interest\textquotedblright{}.

\begin{quote}
\bn{এই} \bn{কাজে} \bn{তার} \bn{কোনো} \bn{স্বার্থ} \bn{নেই।}
\end{quote}

\begin{quote}
\emph{ei kāje tār kono svārthô nei.}
\end{quote}

\begin{quote}
He has no personal interest in this work.
\end{quote}

\begin{cousinweb}
\textbf{\bn{স্বার্থ}} (\emph{svārthô}) is Sanskrit \emph{svārtha}, \emph{sva} 'own' with \emph{artha} 'aim, gain': one's own advantage (Monier-Williams).
\end{cousinweb}

\begin{grammarlens}[title={What you've built}]
One more word for giving reasons and opinions.
\end{grammarlens}

\subsection*{Your turn}
\begin{itemize}
\raggedright
  \item \emph{Say it:} \emph{svārthô}
  \item \emph{Say it:} \emph{svārthô}, once more
  \item \emph{Say it:} \emph{ei kāje tār kono svārthô nei.} — He has no personal interest in this work.
  \item \emph{Recall:} say the Bengali for a right, then the Bengali for a responsibility, then say \emph{svārthô} again
\end{itemize}

\begin{tcolorbox}[breakable,colback=teal!4,colframe=teal!35!black,title={Before you move on}]
What does \bn{স্বার্থ} mean? (\textquotedblleft{}Self-interest; one's own interest\textquotedblright{}.) Say it once more.
\end{tcolorbox}
\section[\texorpdfstring{\bn{রক্ষা} \bn{করা} (rôkkhā kôrā)}{রক্ষা করা (rôkkhā kôrā)}]{\bn{রক্ষা} \bn{করা} (rôkkhā kôrā) — to protect; to defend}

\label{lesson:BN-C339-rokkha-kora}



\begin{quote}
Before the new one: say the Bengali for a responsibility, then the Bengali for self-interest.
\end{quote}

\subsection*{You'll want to know: \bn{রক্ষা} \bn{করা}}
\textbf{\bn{রক্ষা} \bn{করা}} — \emph{rôkkhā kôrā} — \textquotedblleft{}to protect; to defend\textquotedblright{}.

\begin{quote}
\bn{আইন} \bn{আমাদের} \bn{অধিকার} \bn{রক্ষা} \bn{করে।}
\end{quote}

\begin{quote}
\emph{āin āmāder ôdhikār rôkkhā kôre.}
\end{quote}

\begin{quote}
The law protects our rights.
\end{quote}

Past (he/she …): \bn{তিনি} \bn{গ্রামটা} \bn{রক্ষা} \bn{করলেন}, \emph{tini grāmṭā rôkkhā kôrlen}.

\begin{cousinweb}
\textbf{\bn{রক্ষা} \bn{করা}} (\emph{rôkkhā kôrā}) is Sanskrit \emph{rakṣā}, guarding or protection (Monier-Williams), with \bn{করা}. Its everyday cousin \bn{রাখা} 'to keep' comes from the same root.
\end{cousinweb}

\begin{grammarlens}[title={What you've built}]
One more word for giving reasons and opinions.
\end{grammarlens}

\subsection*{Your turn}
\begin{itemize}
\raggedright
  \item \emph{Say it:} \emph{rôkkhā kôrā}
  \item \emph{Say it:} \emph{rôkkhā kôrā}, once more
  \item \emph{Say it:} \emph{āin āmāder ôdhikār rôkkhā kôre.} — The law protects our rights.
  \item \emph{Recall:} say the Bengali for a responsibility, then the Bengali for self-interest, then say \emph{rôkkhā kôrā} again
\end{itemize}

\begin{tcolorbox}[breakable,colback=teal!4,colframe=teal!35!black,title={Before you move on}]
What does \bn{রক্ষা} \bn{করা} mean? (\textquotedblleft{}To protect; to defend\textquotedblright{}.) Say it once more.
\end{tcolorbox}
\section[\texorpdfstring{\bn{ভোগ} \bn{করা} (bhog kôrā)}{ভোগ করা (bhog kôrā)}]{\bn{ভোগ} \bn{করা} (bhog kôrā) — to enjoy a right; to undergo (hardship)}

\label{lesson:BN-C339-bhog-kora}



\begin{quote}
Before the new one: say the Bengali for self-interest, then the Bengali for to defend.
\end{quote}

\subsection*{You'll want to know: \bn{ভোগ} \bn{করা}}
\textbf{\bn{ভোগ} \bn{করা}} — \emph{bhog kôrā} — \textquotedblleft{}to enjoy a right; to undergo (hardship)\textquotedblright{}.

\begin{quote}
\bn{সব} \bn{নাগরিক} \bn{সমান} \bn{অধিকার} \bn{ভোগ} \bn{করে।}
\end{quote}

\begin{quote}
\emph{sôb nāgôrik sômān ôdhikār bhog kôre.}
\end{quote}

\begin{quote}
All citizens enjoy equal rights.
\end{quote}

Past (he/she …): \bn{তিনি} \bn{অনেক} \bn{সুবিধা} \bn{ভোগ} \bn{করলেন}, \emph{tini ônek subidhā bhog kôrlen}.

\begin{cousinweb}
\textbf{\bn{ভোগ} \bn{করা}} (\emph{bhog kôrā}) is Sanskrit \emph{bhoga}, enjoyment or the use of a thing (Monier-Williams), with \bn{করা}. It is the formal verb for having the use of rights and benefits.
\end{cousinweb}

\begin{grammarlens}[title={What you've built}]
One more word for giving reasons and opinions.
\end{grammarlens}

\subsection*{Your turn}
\begin{itemize}
\raggedright
  \item \emph{Say it:} \emph{bhog kôrā}
  \item \emph{Say it:} \emph{bhog kôrā}, once more
  \item \emph{Say it:} \emph{sôb nāgôrik sômān ôdhikār bhog kôre.} — All citizens enjoy equal rights.
  \item \emph{Recall:} say the Bengali for self-interest, then the Bengali for to defend, then say \emph{bhog kôrā} again
\end{itemize}

\begin{tcolorbox}[breakable,colback=teal!4,colframe=teal!35!black,title={Before you move on}]
What does \bn{ভোগ} \bn{করা} mean? (\textquotedblleft{}To enjoy a right; to undergo (hardship)\textquotedblright{}.) Say it once more.
\end{tcolorbox}
